% ---------------------------------------------------------------
\section{Data}
\label{sec:data}
% ---------------------------------------------------------------

\subsection{Source: OpenAlex}

All bibliometric data are drawn from OpenAlex, a fully open index of
scholarly literature maintained by OurResearch~\citep{priem2022openalex}.
OpenAlex aggregates metadata for over 260 million scholarly works, including
authors, source journals, institutional affiliations, topic classifications,
and citation links.  Its coverage spans all scientific fields and extends from
the early twentieth century to the present.  We use a parquet snapshot taken
in early 2025, comprising works, authorships, sources, institutions, topic
assignments, and citation reference lists.  The snapshot is stored as a
collection of columnar parquet files accessed via the DuckDB analytical query
engine.

OpenAlex assigns each work a set of topic tags from a three-level hierarchy:
\emph{topics} are the finest level, grouped into \emph{subfields}, which are
grouped into 26 \emph{fields}.  Each topic tag carries a numeric score
reflecting the strength of the thematic signal in the work; scores are
non-negative and sum across all topics assigned to a given work.  The 26 fields
form the primary grouping level for this study; the field taxonomy and
index values used throughout are listed in Table~\ref{tab:fields}.

\begin{table}[ht]
\centering
\small
\caption{OpenAlex field taxonomy used in this study ($f_{idx}$ = field index).}
\label{tab:fields}
\begin{tabular}{rll}
\toprule
$f_{idx}$ & Field & OA Domain \\
\midrule
11 & Agricultural and Biological Sciences   & Life Sciences \\
12 & Arts and Humanities                    & Social Sciences \\
13 & Biochemistry, Genetics and Molecular Biology & Life Sciences \\
14 & Business, Management and Accounting    & Social Sciences \\
15 & Chemical Engineering                   & Physical Sciences \\
16 & Chemistry                              & Physical Sciences \\
17 & Computer Science                       & Physical Sciences \\
18 & Decision Sciences                      & Social Sciences \\
19 & Earth and Planetary Sciences           & Physical Sciences \\
20 & Economics, Econometrics and Finance    & Social Sciences \\
21 & Energy                                 & Physical Sciences \\
22 & Engineering                            & Physical Sciences \\
23 & Environmental Science                  & Physical Sciences \\
24 & Immunology and Microbiology            & Life Sciences \\
25 & Materials Science                      & Physical Sciences \\
26 & Mathematics                            & Physical Sciences \\
27 & Medicine                               & Health Sciences \\
28 & Neuroscience                           & Life Sciences \\
29 & Nursing                                & Health Sciences \\
30 & Pharmacology, Toxicology and Pharmaceutics & Life Sciences \\
31 & Physics and Astronomy                  & Physical Sciences \\
32 & Psychology                             & Social Sciences \\
33 & Social Sciences                        & Social Sciences \\
34 & Veterinary                             & Health Sciences \\
35 & Dentistry                              & Health Sciences \\
36 & Health Professions                     & Health Sciences \\
\bottomrule
\end{tabular}
\end{table}

\subsection{Work and source filters}

Not all documents in OpenAlex are admitted to the corpus.  We retain only
\emph{articles} and \emph{reviews} (the OpenAlex \texttt{type} field), thereby
excluding preprints, book chapters, data sets, dissertations, and other
document types whose citation behaviour differs systematically from peer-reviewed
journal literature.  Works flagged as paratext (editorials, front matter, indices)
or as retracted are excluded, as are works with zero recorded references, since
a citer with no outgoing edges contributes no signal to the ranking.

The source of publication must be a \emph{journal}, \emph{conference}, or
\emph{book series} in the OpenAlex typology.  This admits the principal channels
of peer-reviewed scientific output while excluding repositories, encyclopaedias,
and trade publications.

\subsection{Institutional authorship and weighting}

OpenAlex resolves each authorship record to an author entity and, where
possible, to one or more institutional affiliations.  Only authorships with a
resolved institution are admitted.  Institutions are further filtered by type:
we retain \emph{education}, \emph{nonprofit}, \emph{government},
\emph{healthcare}, and \emph{other} organisations, and exclude
\emph{company}, \emph{funder}, and \emph{archive} types.  The inclusion of
healthcare institutions captures major clinical research centres --- such as
major cancer centres, hospital systems, and research institutes --- which
OpenAlex classifies separately from universities and which carry substantial
scientific output in biomedical fields.  Company-affiliated works are excluded
because commercial research operates under disclosure constraints that make the
citation environment non-comparable with the academic sector.

A work must have at least one retained institutional authorship to be included;
works whose entire authorship is unresolved, or whose resolved institutions are
all of excluded types, are dropped.

Two institution-weight conventions are computed for each (work, institution)
pair and both are stored in the flat works table.  The \emph{author-fractional}
weight distributes credit for each work across institutions in proportion to
the author share: if work $p$ has $n_a$ authors and author $i$ is affiliated
with $k_i$ retained institutions, then institution $u$'s weight is
\[
  w_{u}^{\mathrm{af}} \;=\; \sum_{i\,:\,u \in U_i} \frac{1}{n_a\, k_i},
\]
so that $\sum_u w_u^{\mathrm{af}} = 1$ per work.  The \emph{direct} weight
distributes credit uniformly across the $K$ distinct retained institutions
affiliated with the work: $w_u^{\mathrm{dir}} = 1/K$.  The author-fractional
weight ($\omega = 0$) is the baseline; the direct weight ($\omega = 1$) is
available as a sensitivity variant.

\subsection{Field weights}

A work's topic scores are aggregated to the field level by summing scores
across all topics belonging to each field, then normalising by the total score
across all topics of the work.  The resulting field weight $f_{p,f} \in [0,1]$
satisfies $\sum_f f_{p,f} = 1$ per work: a work appearing primarily in one
field carries $f_{p,f} \approx 1$ for that field, while an interdisciplinary
work distributes its weight across multiple fields.  This fractional treatment
avoids double-counting citation mass across fields: a work that straddles
Medicine and Neuroscience, for example, contributes less than one full unit of
activity to each field's corpus rather than one full unit to both.

\subsection{Flat works table}

The above filters and weight computations are applied in a single multi-pass
DuckDB query over the raw parquet snapshot, producing a \emph{flat works}
table at the granularity of one row per (work $\times$ institution $\times$
subfield).  This granularity captures the full joint distribution of
institutional authorship and topical classification without requiring repeated
scans of the raw data at later pipeline stages.  For a work with $K$
institutions and $F$ distinct subfield tags, the flat works table contains
$K \times F$ rows for that work.

The flat works table covers works published between 2000 and 2025, yielding
153.5~million rows.  Each row carries: work identifier, publication year,
source identifier, institution identifier, country code (ISO 3166-1), both
institution weights, subfield identifier and name, field identifier and weight,
the CWTS Leiden main field grouping (one of five groups), and the number of
references in the work.

\subsection{Corpus references}

Citation links are extracted from the OpenAlex references parquet
(${\approx}500$~million rows), which records for each work the set of works it
cites.  We materialise a \emph{corpus references} table containing only those
citation pairs $(p, q)$ where both the citer $p$ and the cited work $q$ appear
in flat works.  This within-corpus filter retains approximately 358~million
citation pairs and is computed once and reused across all 26 fields, avoiding
repeated full scans of the raw reference data.

\subsection{Census window and candidacy}

The ranking is computed over a five-year census window, 2020--2024.  Only
works published within this window enter the corpus as citers; cited works,
however, may have been published at any time provided they appear in flat works
(i.e.\ they satisfy the type, source, and institution filters).  In practice
the cited works window is 2016--2024, covering papers published up to nine
years before the end of the census window and still generating citations within
it.

For each field $f$ and census window $[t_0, t_1]$, the candidacy of a source
$s$ is its \emph{weighted works}: the sum of field weights $f_{p,f}$ across
all works $p$ with $\mathrm{source}(p) = s$ and $t_0 \leq
\mathrm{year}(p) \leq t_1$.  The source candidacy computation deduplicates
at the (work, source) level before summing, because flat works has multiple
rows per work (one per institution), and naïve summation would otherwise count
each work's field weight once per affiliated institution.

The institution candidacy is the sum of $f_{p,f} \cdot w_u^{\mathrm{af}}$
across all (work, institution, subfield) rows for institution $u$ in the
window: each work's field-weight contribution is fractionally allocated to each
affiliated institution, so no deduplication is required.

Candidacy tables are precomputed once per window for all 26 fields and reused
across all ranking runs, including sensitivity variants and country-bloc runs.

\subsection{Scale}

Table~\ref{tab:scale} reports the number of retained sources and institutions
after the $\tau_s = \tau_u = 10$ per year threshold for each of the five CWTS
Leiden main fields under the 2020--2024 window.  These figures convey the
scale of the ranking problem: the bipartite matrix $M_S = H_{SI}H_{IS}$ is
$N_s \times N_s$ dense in its sparsity pattern, so field sizes in the range of
$10^3$--$10^4$ are computationally tractable while still providing
statistically reliable rankings.

\begin{table}[ht]
\centering
\caption{Retained sources ($N_s$) and institutions ($N_u$) per CWTS Leiden main field,
         window 2020--2024, $\tau = 10$ per year.}
\label{tab:scale}
\begin{tabular}{rlrr}
\toprule
Group & Field & $N_s$ & $N_u$ \\
\midrule
1 & Mathematics and Computer Science  &  3,624 &  3,754 \\
2 & Physical Sciences and Engineering &  6,724 &  5,694 \\
3 & Life and Earth Sciences           &  8,475 &  7,008 \\
4 & Biomedical and Health Sciences    & 11,585 & 10,069 \\
5 & Social Sciences and Humanities    & 16,082 &  7,251 \\
\bottomrule
\end{tabular}
\end{table}

The individual 26 OA fields are smaller subsets of these Leiden groups and
vary considerably in size: major clinical and life science fields (Medicine,
field 27; Biochemistry \& Genetics, field 13) retain several thousand sources
and institutions each, while smaller fields (Dentistry, field 35; Veterinary,
field 34) retain a few hundred.  Across all 26 fields the ranking procedure
was run independently, with each field receiving its own bipartite citation
graph, candidacy threshold, and Perron eigenvector computation.
